%% ma: L12-B14-0102
%% lop: 12
%% chuong: 5
%% bai: 14
%% dang: 12.14.30
%% loai: TLN
%% muc: VD
%% dap_an: "4"
%% nguon: "Ảnh giáo viên gửi 07/11/2026 (Tư Duy Mở - tuduymo.com, đề tự luyện Phương trình mặt phẳng, đề lập phương trình mặt phẳng), Câu 20"
%% kien_thuc: "Bài 14 (phương trình mặt phẳng), Bài 23-25 lớp 11 (đường thẳng vuông góc với mặt phẳng, tứ diện vuông)"
%% trang_thai: cho-duyet
%% ly_do: "Dạng 12.14.30 mới đề xuất, chờ giáo viên duyệt."
%% ghi_chu: "Tài liệu gốc không kèm đáp án; đáp án do Claude giải và kiểm bằng sympy."
\begin{ex}
Trong không gian $Oxyz$, cho mặt phẳng $(P)\colon ax+by+cz+28=0$ là mặt phẳng cắt ba trục tọa độ $Ox$, $Oy$, $Oz$ lần lượt tại ba điểm $A$, $B$, $C$ sao cho điểm $M(2;3;-1)$ là trực tâm của tam giác $ABC$. Hãy xác định giá trị của biểu thức $T=a-b+c$ (làm tròn kết quả đến hàng đơn vị).
\shortans{4}
\loigiai{Tứ diện $OABC$ có ba cạnh tại $O$ đôi một vuông góc, $M$ là trực tâm tam giác $ABC$ nên $OM\perp(ABC)$.
$(P)$ có vectơ pháp tuyến cùng phương $\overrightarrow{OM}=(2;3;-1)$: $(a;b;c)=t(2;3;-1)$. $M\in(P)\Rightarrow 14t+28=0\Rightarrow t=-2$.
$a=-4$, $b=-6$, $c=2$, $T=-4+6+2=4$.}
\end{ex}
